% Source: https://artofproblemsolving.com/community/c3h2483605_nt_problems
% Stable identity: aops-topic-2483605-q04
Since $k(k+2)=(k+1)^2-1$, the $k$-th term is $(k+1)^4-(k+1)^2$. With $j=k+1$ and $m=n+1$, and since the $j=1$ term $1-1$ is zero,
\[S_n=\sum_{j=1}^{m}\left(j^4-j^2\right).\]
Using
\[\sum_{j=1}^{m}j^2=\dfrac{m(m+1)(2m+1)}{6},\qquad\sum_{j=1}^{m}j^4=\dfrac{m(m+1)(2m+1)\left(3m^2+3m-1\right)}{30},\]
we get
\[S_n=\dfrac{m(m+1)(2m+1)\left(3m^2+3m-1-5\right)}{30}=\dfrac{m(m+1)(2m+1)(m-1)(m+2)}{10}.\]
Substituting $m=n+1$:
\[S_n=\boxed{\dfrac{n(n+1)(n+2)(n+3)(2n+3)}{10}}.\]
Check: $S_1=12=\dfrac{1\cdot2\cdot3\cdot4\cdot5}{10}$ and $S_2=12+72=84=\dfrac{2\cdot3\cdot4\cdot5\cdot7}{10}$.
